\documentclass[conference]{IEEEtran}

\usepackage{iftex}
\ifPDFTeX
\usepackage[utf8]{inputenc}
\usepackage[T1]{fontenc}
\fi
\usepackage{amsmath,amssymb,bm,mathtools}
\usepackage{newtxtext,newtxmath}
\usepackage{textcomp}
\usepackage{array,booktabs,tabularx}
\usepackage{graphicx}
\usepackage{siunitx}
\usepackage{cite}
\usepackage[breaklinks=true,hidelinks]{hyperref}
\usepackage{bookmark}
\usepackage[nameinlink,capitalise,noabbrev]{cleveref}
\usepackage[final]{microtype}

\interdisplaylinepenalty=2500
\allowdisplaybreaks
\emergencystretch=2em
\sisetup{detect-weight=true,detect-family=true,per-mode=symbol}

\DeclareRobustCommand{\vct}[1]{\bm{#1}}
\DeclareRobustCommand{\mat}[1]{\bm{#1}}
\newcommand{\R}{\mathbb{R}}
\newcommand{\T}{^{\mathsf T}}
\newcommand{\Skew}[1]{\left[#1\right]_{\times}}
\newcommand{\norm}[1]{\left\lVert#1\right\rVert}
\newcommand{\diag}{\operatorname{diag}}
\newcommand{\N}{\mathcal N}

\title{GNSS-Aided Fusion in Marine Inertial Navigation Filters\\
\large Absolute Navigation, Lever Arms, Asynchronous Aiding, Outages, and Wave--Trajectory Separation (draft)}

\author{%
\IEEEauthorblockN{Mikhail S. Grushinskiy}
\IEEEauthorblockA{Independent Researcher, USA, Fair Lawn\\
Bareboat Necessities GitHub Project\\
Email: mgrouch@users.github.com}%
}

\begin{document}
\maketitle

\begin{abstract}
Marine inertial motion estimators and absolute navigation filters solve related but
not identical problems.  A wave-motion estimator seeks bounded surge, sway, and
heave about a slowly moving vessel trajectory, whereas GNSS reports the absolute
motion of an antenna phase center in an earth-fixed or local-level frame.  Treating
those quantities as one position state creates a structural conflict: virtual
zero-mean constraints that are useful for wave-drift control are false for a vessel
under way, while direct absolute GNSS updates can suppress or absorb the wave motion
one wishes to estimate.  This paper develops a common GNSS-aiding architecture for
four marine inertial filters in the Bareboat Necessities \texttt{ocean-imu}
repository: the OU--II and OU--III quaternion MEKFs, the right-invariant two-frame
Lie-group EKF (TFG), and the time-varying-gain nonlinear observer (NLO).  The central
proposal is to separate a slow absolute navigation trajectory
$(\vct r_n,\vct u_n,\vct a_n)$ from bounded wave motion
$(\vct p_w,\vct v_w,\vct a_w)$ and let GNSS observe their physical sum at the
antenna.  Known IMU-to-center-of-gravity and GNSS-antenna lever arms are modeled with
full rigid-body kinematics.  The paper also specifies asynchronous event-time fusion,
delayed-measurement handling, outage/reacquisition logic, local-tangent-plane
navigation, and validation protocols.  OU--II and OU--III already contain the IMU
lever-arm acceleration correction and generic full-covariance position/velocity
updates needed as building blocks; the complete absolute-navigation layer and the
corresponding TFG/NLO extensions remain prospective.  No new GNSS performance result
is claimed here; the manuscript defines the model and the experiments required to
establish one.
\end{abstract}

\begin{IEEEkeywords}
GNSS/INS, marine inertial navigation, wave motion, lever arm, asynchronous sensor
fusion, outage reacquisition, multiplicative extended Kalman filter, Lie-group EKF,
nonlinear observer, heave, surge, sway
\end{IEEEkeywords}

\section{Introduction}
\label{sec:introduction}

GNSS and inertial sensors are naturally complementary.  GNSS bounds long-term
position and velocity drift but arrives slowly, with latency, dropouts, multipath,
and a measurement point displaced from the IMU.  An IMU provides high-rate angular
rate and specific force but cannot maintain absolute position without external
aiding.  This complementarity motivates GNSS/INS integration in conventional
navigation filters and in nonlinear observers designed for satellite-aided marine
systems \cite{Kalman1960_Filtering,BryneFossenJohansen2014_TVG_NLO}.

Wave-motion estimation adds a second requirement.  The estimator must preserve
oscillatory motion that an ordinary navigation filter may regard as disturbance.
Marine heave and wave estimators therefore introduce band-limited motion models,
virtual zero constraints, or other drift-suppression structure
\cite{Auestad2013_HeaveStrapdownIMU,Pascoal2009_KalmanDirectionalSpectrum,
YurovskyKudinov2025_IMUwaves}.  Those constraints are meaningful only for a
\emph{wave displacement about a moving reference}, not for the absolute position of
a vessel sailing through the world.

The repository currently contains four relevant estimator families.  OU--II carries
velocity, displacement, and OU acceleration and periodically applies zero-position
and zero-velocity virtual measurements.  OU--III replaces those direct anchors by a
zero pseudo-measurement on the integral of displacement.  TFG expresses the OU--III
world vectors on a two-frame Lie-group error state.  NLO is derived from a
satellite-aided nonlinear observer and already has a horizontal GNSS position path,
with virtual horizontal aiding available when the position reference disappears.
The four filters therefore provide an unusually useful test bed for separating the
\emph{geometry of GNSS aiding} from the \emph{internal drift-control mechanism}.

This paper has six goals.  First, it fixes reference-point semantics so that all
filters estimate motion of the vessel center of gravity (CoG) while sensors may be
mounted elsewhere.  Second, it gives the rigid-body IMU and GNSS antenna lever-arm
models.  Third, it introduces separate absolute-navigation and wave-motion states.
Fourth, it specifies event-time GNSS fusion, including delayed observations and
reacquisition after outages.  Fifth, it maps the architecture onto OU--II, OU--III,
TFG, and NLO without pretending their error coordinates are interchangeable.  Sixth,
it predeclares the studies needed before any GNSS-aided performance claim should be
made.

\section{Frames, Reference Points, and State Semantics}
\label{sec:frames}

Let $\mathcal W$ be a local NED navigation frame, $\mathcal B$ the physical rigid-body
frame of the vessel, and $\mat R_{BW}$ the rotation from body to world.  Its inverse is
$\mat R_{WB}=\mat R_{BW}\T$.  Three physical points are distinguished:
\begin{itemize}
  \item $C$: vessel center of gravity or another declared navigation reference point;
  \item $I$: IMU sensing point;
  \item $G$: GNSS antenna phase center.
\end{itemize}
Their fixed body-frame lever arms are
\begin{equation}
\vct \ell_I^{\mathcal B}=\vct r_I^{\mathcal B}-\vct r_C^{\mathcal B},
\qquad
\vct \ell_G^{\mathcal B}=\vct r_G^{\mathcal B}-\vct r_C^{\mathcal B}.
\label{eq:lever-defs}
\end{equation}
If a survey provides only the IMU-to-antenna baseline,
$\vct \ell_{GI}^{\mathcal B}$, then
\begin{equation}
\vct \ell_G^{\mathcal B}=\vct \ell_I^{\mathcal B}+\vct \ell_{GI}^{\mathcal B}.
\label{eq:lever-compose}
\end{equation}
The sign convention in \cref{eq:lever-defs} is worth fixing explicitly because a
reversed baseline produces a plausible but systematically wrong rotational
correction.

The core wave states in the existing OU filters should be interpreted as motion of
$C$ about a moving reference.  They are not geodetic coordinates.  Introduce a slow
navigation trajectory
\begin{equation}
\vct x_n=
\begin{bmatrix}
\vct r_n\T & \vct u_n\T & \vct a_n\T
\end{bmatrix}\T,
\label{eq:nav-state}
\end{equation}
and bounded wave motion
\begin{equation}
\vct x_w=
\begin{bmatrix}
\vct p_w\T & \vct v_w\T & \vct a_w\T
\end{bmatrix}\T,
\label{eq:wave-state}
\end{equation}
with the OU--III extension $\vct S_w=\int\vct p_w\,dt$.  The physical CoG motion is
then
\begin{align}
\vct r_C^{\mathcal W} &= \vct r_n^{\mathcal W}+\vct p_w^{\mathcal W},
\label{eq:sum-position}\\
\vct v_C^{\mathcal W} &= \vct u_n^{\mathcal W}+\vct v_w^{\mathcal W},
\label{eq:sum-velocity}\\
\vct a_C^{\mathcal W} &= \vct a_n^{\mathcal W}+\vct a_w^{\mathcal W}.
\label{eq:sum-acceleration}
\end{align}
Equations~\eqref{eq:sum-position}--\eqref{eq:sum-acceleration} are the organizing
principle of the paper.  Absolute navigation and wave motion are separate latent
components whose sum is physical.

\section{Rigid-Body Lever-Arm Kinematics}
\label{sec:rigid-body}

For any fixed point $P$ on the rigid body with body-frame lever arm
$\vct\ell_P^{\mathcal B}$ from $C$, the velocity and acceleration are
\begin{align}
\vct v_P^{\mathcal W}
&=\vct v_C^{\mathcal W}
 +\mat R_{BW}\bigl(\vct\omega^{\mathcal B}\times
 \vct\ell_P^{\mathcal B}\bigr),
\label{eq:lever-velocity}\\
\vct a_P^{\mathcal W}
&=\vct a_C^{\mathcal W}
 +\mat R_{BW}\Bigl(
 \vct\alpha^{\mathcal B}\times\vct\ell_P^{\mathcal B}
 +\vct\omega^{\mathcal B}\times
  (\vct\omega^{\mathcal B}\times\vct\ell_P^{\mathcal B})
 \Bigr).
\label{eq:lever-acceleration}
\end{align}
The first rotational term is tangential acceleration and the second is centripetal.
There is no Coriolis term because the sensor is fixed to the rigid hull.

Two practical details matter.  Angular acceleration $\vct\alpha$ is usually obtained
by differentiating a bias-corrected gyro signal, which amplifies noise.  It therefore
needs timestamp-consistent differentiation and often mild smoothing.  Second,
lever-arm corrections must use the same physical-body convention as the surveyed
baseline.  If an estimator internally uses a virtual un-heeled body frame, the lever
arm should be transformed into that frame before applying the kinematics; the
surveyed geometry itself should remain stored in the physical body frame.

\section{IMU Lever Arm from the Center of Gravity}
\label{sec:imu-lever}

An accelerometer mounted at $I$ measures specific force at the IMU point, not at the
CoG.  Ignoring scale-factor and misalignment errors, its model is
\begin{align}
\vct z_f^{\mathcal B}
={}&\mat R_{WB}
\bigl(\vct a_C^{\mathcal W}-\vct g^{\mathcal W}\bigr)
+\vct a_{I/C}^{\mathcal B}
+\vct b_a+\vct n_a,
\label{eq:imu-specific-force}\\
\vct a_{I/C}^{\mathcal B}
={}&\vct\alpha^{\mathcal B}\times\vct\ell_I^{\mathcal B}
+\vct\omega^{\mathcal B}\times
 (\vct\omega^{\mathcal B}\times\vct\ell_I^{\mathcal B}).
\label{eq:imu-rot-acc}
\end{align}
Thus a filter whose translational state is defined at the CoG should either subtract
$\vct a_{I/C}^{\mathcal B}$ from the measurement before fusion or add it to the
predicted IMU specific force.  The latter is preferable inside a coupled MEKF because
its dependence on attitude and gyro bias can be represented in the measurement
Jacobian.

\subsection{Current OU--II and OU--III implementation}

The current OU--II and OU--III classes already implement this model.  Both expose
\texttt{set\_imu\_lever\_arm\_body()},
\texttt{clear\_imu\_lever\_arm()}, and
\texttt{set\_alpha\_smoothing\_tau()}.  The configured vector is documented as the
IMU position with respect to CoG in the physical body frame.  During the accelerometer
measurement update the implementation forms
\begin{equation}
\vct a_{I/C}
=\vct\alpha\times\vct\ell_I
 +\vct\omega\times(\vct\omega\times\vct\ell_I),
\label{eq:repo-imu-lever}
\end{equation}
from bias-corrected angular rate and the internal angular-acceleration estimate.  When
the optional wind-heel frame is active, the physical lever arm is de-heeled into the
virtual frame before the term is evaluated.  The accelerometer Jacobian also includes
the gyro-bias sensitivity of the lever-arm term.

This means the first GNSS-aided implementation should \emph{not} add a second,
independent IMU lever-arm correction in a wrapper.  Doing so would double-correct the
measurement.  Instead, the GNSS layer should keep the OU filter's existing CoG
specific-force model and add only the antenna-side geometry in the GNSS observation.

\subsection{Limits of the current correction}

The lever arm should be disabled during initialization if the angular-acceleration
estimate is not yet meaningful; the current warm-up path already clears it.  A
production study should quantify the trade between angular-acceleration smoothing and
phase error.  It should also test lever-arm miscalibration and a slowly shifting CoG
due to crew, fuel, or payload.  Online lever-arm estimation is possible in principle,
but it requires rotational excitation and creates additional observability coupling;
it is outside the first implementation proposed here.

\section{GNSS Antenna Measurement Model}
\label{sec:gnss-measurement}

Assume a GNSS position solution has been converted to the local frame and refers to
the antenna phase center.  Using \cref{eq:sum-position}, the position measurement is
\begin{equation}
\boxed{
\vct z_r
=\vct r_n+\vct p_w
 +\mat R_{BW}\vct\ell_G^{\mathcal B}
 +\vct n_r.}
\label{eq:gnss-position}
\end{equation}
If the receiver supplies Doppler velocity at the antenna,
\begin{equation}
\boxed{
\vct z_v
=\vct u_n+\vct v_w
 +\mat R_{BW}
   (\vct\omega^{\mathcal B}\times\vct\ell_G^{\mathcal B})
 +\vct n_v.}
\label{eq:gnss-velocity}
\end{equation}
The position lever arm is attitude dependent; the velocity lever arm is attitude,
angular-rate, and gyro-bias dependent.  During a turn, omitting the latter can create
a velocity residual that looks exactly like sway or wave orbital velocity.

For a multiplicative EKF, the position Jacobian contains identity blocks for both
$\vct r_n$ and $\vct p_w$, plus an attitude block proportional to the skew matrix of
the rotated lever arm.  The velocity Jacobian similarly contains identity blocks for
$\vct u_n$ and $\vct v_w$ and an attitude block generated by
$\vct\omega\times\vct\ell_G$.  If $\vct\omega$ is formed from gyro measurement minus
estimated bias, there is also a gyro-bias block.  The exact signs depend on the
filter's multiplicative-error convention; they must be derived in each filter's own
coordinates rather than copied between OU and TFG implementations.

The antenna acceleration, useful for diagnostics and latency propagation, follows
\cref{eq:lever-acceleration} with $P=G$:
\begin{equation}
\vct a_G^{\mathcal W}=\vct a_C^{\mathcal W}
+\mat R_{BW}\Bigl(
\vct\alpha\times\vct\ell_G
+\vct\omega\times(\vct\omega\times\vct\ell_G)
\Bigr).
\label{eq:gnss-antenna-acc}
\end{equation}

\section{Absolute Navigation in a Local Tangent Plane}
\label{sec:absolute-navigation}

GNSS fixes should be converted from latitude, longitude, and height to a declared
navigation frame before filtering.  For local marine operation a fixed NED origin is
usually sufficient:
\begin{equation}
\vct z_r^{\mathcal W}
=\operatorname{LLAtoNED}
\bigl(\operatorname{LLA}_{G};\operatorname{LLA}_0\bigr).
\label{eq:lla-ned}
\end{equation}
The filter must record whether the vertical coordinate is ellipsoidal height,
orthometric height, or a local mean-water reference.  Mixing those conventions can
look like a slow heave bias or tide signal.

For long routes, a single flat local tangent plane eventually becomes inaccurate.
Two clean options are an ECEF navigation state or periodic local-frame re-anchoring.
Re-anchoring is a coordinate transformation, not a measurement reset: the navigation
position, velocity, covariance, and any world-frame cross-covariances must be rotated
and translated consistently.  The bounded wave state is translation invariant and
should not jump when the geodetic origin changes.

Receiver-reported horizontal/vertical accuracy can initialize
$\mat R_r$ and Doppler-velocity accuracy can initialize $\mat R_v$.  DOP values alone
are not measurement covariance unless converted with an assumed ranging-error model.
The first implementation should support anisotropic NED covariances and allow vertical
GNSS aiding to be disabled independently.

\section{Joint Separation of Vessel Trajectory and Wave Motion}
\label{sec:joint-separation}

The preferred architecture is a single joint estimator in which navigation and wave
states share covariance or observer coupling.  A useful slow trajectory model is
\begin{align}
\dot{\vct r}_n &= \vct u_n,\\
\dot{\vct u}_n &= \vct a_n,\\
d\vct a_n &= -\mat\Lambda_n\vct a_n\,dt+\mat L_n\,d\vct W_n,
\label{eq:nav-process}
\end{align}
where $\mat\Lambda_n$ is chosen much slower than the wave-acceleration dynamics.  A
random-walk acceleration model is another option.  The wave channel keeps the filter's
existing OU or integrated-wave process.

The separation is therefore model based rather than produced by a post-hoc low-pass
filter.  GNSS position and velocity observe the sums in
\cref{eq:gnss-position,eq:gnss-velocity}; the IMU observes the acceleration sum in
\cref{eq:sum-acceleration}.  Cross-covariance lets GNSS constrain accelerometer bias,
low-frequency attitude leakage, and trajectory acceleration without forcing the
bounded wave state to carry secular motion.

The split is not universally identifiable.  A sharp vessel maneuver can occupy the
same frequency band as a wave group, and a long swell can overlap a slowly varying
trajectory model.  Separation then depends on additional structure: GNSS velocity,
heading, water speed, control inputs, current models, the wave OU prior, and the
relative process covariances.  The ambiguity should be measured, not hidden by a
fixed cutoff.

\subsection{Cascaded alternative}

A lower-complexity implementation can run a slow GNSS navigation filter first,
estimate $(\vct r_n,\vct u_n,\vct a_n)$, subtract the estimated trajectory acceleration
from the leveled IMU signal, and feed the residual to an unchanged wave filter.  This
is attractive for initial prototyping, but it ignores cross-covariance and can imprint
the navigation filter's phase lag into the wave estimate.  The cascaded architecture
should therefore be treated as an ablation or engineering approximation, not as
mathematically equivalent to the joint filter.

\section{Asynchronous GNSS Aiding and Timestamp Alignment}
\label{sec:asynchronous}

The IMU may run near hundreds of hertz while GNSS arrives at 1--20~Hz.  A discrete
Kalman implementation should fuse a GNSS sample exactly once at its measurement time.
Holding the same fix and reapplying it at every IMU tick double-counts information.
The preferred event sequence is
\begin{enumerate}
  \item propagate IMU states to the GNSS measurement timestamp $t_g$;
  \item evaluate attitude, angular rate, and both lever arms at $t_g$;
  \item apply the GNSS position/velocity update once;
  \item continue propagation to the current processing time.
\end{enumerate}

If the receiver delivers delayed observations, a short state-history ring buffer can
support rewind/update/replay.  For constrained embedded systems, a cheaper
approximation is to propagate the GNSS measurement to the current epoch using its
velocity and measured age, then inflate covariance for acceleration and timestamp
uncertainty.  The approximation should be tested against exact replay.

A clock offset $\delta t$ is especially damaging under motion.  To first order,
position error contains $\vct v_G\delta t$ and velocity error contains
$\vct a_G\delta t$.  At \SI{5}{m/s}, a \SI{100}{ms} timing error already creates a
\SI{0.5}{m} apparent position offset before any lever-arm effect is considered.

NLO requires a slightly different interpretation.  Its current satellite reference
is a continuously evaluated observer input rather than a discrete Kalman update.
Sample-and-hold GNSS can therefore be legitimate, but the implementation still needs
sample age, freshness timeout, and explicit invalidation during outages.  A stale
position reference must not remain ``valid'' indefinitely.

\section{Application to OU--II}
\label{sec:ou2}

The current OU--II translational state consists of wave velocity, wave displacement,
and latent OU acceleration, with accelerometer bias optional.  Periodically, the
implementation calls the same full-covariance position and velocity update routines
with zero measurements.  That architecture makes the semantic boundary very clear:
those states are bounded wave states.

A GNSS-aided OU--II should therefore augment, rather than reinterpret, the core state:
\begin{equation}
\vct x_{\rm OU2+nav}=
\begin{bmatrix}
\delta\vct\theta & \vct b_g &
\vct v_w & \vct p_w & \vct a_w & \vct b_a &
\vct u_n & \vct r_n & \vct a_n
\end{bmatrix}\T.
\label{eq:ou2-augmented}
\end{equation}
The existing zero-$p$ and zero-$v$ virtual measurements remain attached only to
$(\vct p_w,\vct v_w)$.  They must never be applied to $(\vct r_n,\vct u_n)$.

The existing
\texttt{measurement\_update\_position\_pseudo()} and
\texttt{measurement\_update\_velocity\_pseudo()} routines demonstrate the required
Joseph-form full-covariance update machinery, but a direct call with absolute GNSS
position would still be wrong because those methods currently target the wave blocks.
A true GNSS update needs a new measurement matrix that contains both navigation and
wave blocks and the antenna lever-arm attitude terms from
\cref{eq:gnss-position,eq:gnss-velocity}.

The existing IMU CoG lever-arm model should remain enabled.  In the augmented
accelerometer prediction, the world acceleration entering
\cref{eq:imu-specific-force} becomes $\vct a_n+\vct a_w$.  This prevents a sustained
turn or propulsion transient from being misidentified as wave acceleration.

\section{Application to OU--III}
\label{sec:ou3}

OU--III adds the wave integral state $\vct S_w$ and regulates drift through a virtual
$\vct S_w=0$ measurement rather than direct $p_w=0$ and $v_w=0$ updates.  The GNSS
extension should preserve that distinction:
\begin{equation}
\vct x_{\rm OU3+nav}=
\begin{bmatrix}
\delta\vct\theta & \vct b_g &
\vct v_w & \vct p_w & \vct S_w & \vct a_w & \vct b_a &
\vct u_n & \vct r_n & \vct a_n
\end{bmatrix}\T.
\label{eq:ou3-augmented}
\end{equation}
The integral state is
\begin{equation}
\dot{\vct S}_w=\vct p_w,
\label{eq:sw-wave-only}
\end{equation}
not the integral of absolute position.  Integrating $\vct r_n+\vct p_w$ would make the
zero-$S$ regularizer fundamentally inconsistent with a vessel under way.

GNSS position and velocity constrain the sums
$(\vct r_n+\vct p_w)$ and $(\vct u_n+\vct v_w)$.  The OU--III integral regularizer
continues to bound the wave component through periods of GNSS loss.  Under reliable
GNSS, one may consider softening the $S_w=0$ measurement because GNSS provides an
additional low-frequency anchor, but this is not obviously beneficial: the GNSS
trajectory/wave split and the integral regularizer solve different problems.  The
correct policy should be determined by an ablation comparing fixed $R_S$, GNSS-aware
$R_S$ inflation, and temporary $S$ suppression.

The present OU--III implementation already contains the same CoG IMU lever-arm
compensation as OU--II, including transformation through the optional un-heeled body
frame.  Therefore the new work is chiefly state augmentation, GNSS antenna geometry,
and asynchronous measurement handling rather than a new accelerometer correction.

\section{Application to the Two-Frame Lie-Group EKF}
\label{sec:tfg}

The current TFG estimator stores four world vectors
$(\vct v_w,\vct p_w,\vct S_w,\vct a_w)$ inside a two-frame group and uses a
right-invariant error-state EKF.  It does not currently expose the OU-style generic
GNSS position/velocity measurement API, so the GNSS update must be derived in TFG
coordinates rather than copied from the MEKF.

A natural joint extension is to add three more world-vector columns,
$(\vct u_n,\vct r_n,\vct a_n)$, producing a seven-column two-frame state.  The physical
measurement functions remain exactly
\cref{eq:gnss-position,eq:gnss-velocity}; only the tangent-coordinate Jacobians change.
The GNSS position residual must couple attitude, $\vct p_w$, and $\vct r_n$; the
velocity residual must couple attitude, gyro bias, $\vct v_w$, and $\vct u_n$.

The update should preserve TFG's full retraction/reset covariance transport.  An
additive EKF attitude Jacobian inserted without the corresponding group reset would
be internally inconsistent.  Conversely, adding navigation columns does not justify
a stronger global invariant-filter claim: the repository's complete marine process
already contains estimate-dependent OU/bias couplings and is intentionally described
as an invariant-style right-invariant EKF rather than a globally log-linear InEKF.

A lower-risk first implementation can append a conventional additive navigation block
to the TFG covariance while retaining invariant attitude/wave coordinates.  That
hybrid design should be compared with the seven-column group extension.  The study
should report whether either choice materially changes consistency or merely code
complexity.

\section{Application to the Time-Varying-Gain NLO}
\label{sec:nlo}

The NLO has the strongest historical connection to GNSS because its source formulation
is explicitly an inertial-navigation observer aided by satellite reference systems
\cite{BryneFossenJohansen2014_TVG_NLO}.  The repository implementation currently
exposes a horizontal \texttt{GnssXY} object containing position, RMS, and a validity
flag.  When a usable GNSS reference is absent, it can switch the horizontal axes to
the same virtual zero-mean integrated-position concept used for the vertical channel.
This already provides an interesting outage fallback.

However, the present interface does not yet provide GNSS velocity, vertical absolute
navigation, antenna lever-arm fields, or an explicit GNSS timestamp/freshness model.
Those should be added before comparing NLO fairly with the Kalman families.

The simplest NLO lever-arm integration is a measurement front end.  At the GNSS
timestamp, correct antenna position and velocity to the CoG:
\begin{align}
\vct z_{r,C} &= \vct z_{r,G}-\hat{\mat R}_{BW}\vct\ell_G,\\
\vct z_{v,C} &= \vct z_{v,G}
 -\hat{\mat R}_{BW}(\hat{\vct\omega}\times\vct\ell_G).
\label{eq:nlo-cog-front-end}
\end{align}
Because NLO is not covariance based, uncertainty from attitude, lever-arm calibration,
and angular-rate error cannot simply appear as cross-covariance.  It should instead
inflate the effective reference uncertainty or modulate the observer gain.

Trajectory/wave separation is a deeper change.  The current horizontal position state
with GNSS aiding behaves as a navigation quantity, while the no-GNSS fallback uses a
zero-mean wave-style constraint.  A unified GNSS/wave product should introduce a slow
trajectory state and a bounded wave state rather than asking one state to change
meaning when GNSS validity changes.  This requires re-deriving the translational
observer gains or introducing an explicitly complementary two-channel observer.  It
should not be implemented by merely high-pass filtering the published position output
and calling the remainder ``wave motion.''

\section{GNSS Outages and Reacquisition}
\label{sec:outages}

A robust implementation should expose three aiding modes:
\begin{description}
  \item[AIDED:] fresh GNSS measurements pass quality and innovation gates;
  \item[COAST:] GNSS is absent or rejected; inertial/navigation states propagate and
  uncertainty grows;
  \item[REACQUIRE:] new fixes have returned, but the estimator has not yet accepted a
  stable reference.
\end{description}

During COAST, the bounded wave regularizer should continue normally.  The absolute
navigation state must \emph{not} be pulled toward zero; its covariance or observer
uncertainty should grow according to the navigation process model.  This asymmetry is
one of the main benefits of explicit wave/trajectory separation.

Reacquisition should use receiver quality and innovation consistency.  For a Kalman
filter, normalized innovation squared (NIS) can gate position and velocity separately.
Velocity is often safer to accept before position because a long outage may produce a
large absolute position residual while Doppler velocity is already physically
plausible.  A soft-start policy can inflate $\mat R_r$ and $\mat R_v$ for the first few
fixes, then return to nominal values as consistency is established.

If an outage is so long that the navigation position must be reset, the correction
should be applied to $\vct r_n$ while preserving $\vct p_w$.  Resetting the wave state
would create an artificial phase discontinuity in heave/surge/sway.  Attitude should
not be hard-reset from single-antenna GNSS position.  Course over ground can provide a
heading cue only when speed and sideslip assumptions justify it; dual-antenna GNSS is
a separate heading measurement and should be modeled as such.

\section{Observability and Ambiguities}
\label{sec:observability}

GNSS dramatically improves low-frequency observability, but several ambiguities
remain.  Position alone does not instantly separate $\vct r_n$ from $\vct p_w$ because
only their sum is measured.  The process spectra, velocity measurement, and wave
regularizer create the separation over time.  Likewise, horizontal specific-force
error can be explained by accelerometer bias, tilt error, slow maneuver acceleration,
or wave acceleration.  GNSS velocity is especially valuable because it constrains
that ambiguity one derivative closer to acceleration than position alone.

Lever arms create both information and sensitivity.  A known nonzero antenna lever
arm can couple GNSS residuals to attitude during rotational excitation.  But an
incorrect lever arm produces attitude-correlated position and velocity residuals that
can be absorbed by wave motion or bias states.  The correct question is therefore not
whether lever arms are ``small,'' but whether their induced motion is small relative
to the requested wave/navigation accuracy.

The proposed study should compute finite-horizon empirical observability Gramians or
singular values for at least the following cases: position only; position plus Doppler
velocity; magnetometer enabled/disabled; zero versus nonzero antenna lever arm; steady
straight motion versus turns; and GNSS aided versus outage.  For TFG, the analysis
should be performed in its local invariant error coordinates.  For NLO, use the
linearized observer error dynamics rather than importing a Kalman observability test
unchanged.

\section{Interaction with Wind Heel and Other Reference-Point Transforms}
\label{sec:heel-interaction}

If the optional steady-wind-heel transform is active, both lever arms remain surveyed
in the physical body frame.  OU--II and OU--III already de-heel the IMU lever arm
internally before evaluating rotational acceleration.  The GNSS antenna term must use
the physical boat attitude when predicting the physical antenna position.  A virtual
un-heeled attitude must therefore be recomposed with the current heel estimate before
applying $\mat R_{BW}\vct\ell_G$.

This distinction prevents a subtle double transformation: the IMU measurement model
may be evaluated in the virtual frame, while GNSS geometry belongs to the real hull.
The validation should include static heel plus mast/rail antenna offsets and verify
that changing the estimated heel does not create a fictitious GNSS translation.

\section{Implementation Roadmap}
\label{sec:roadmap}

A staged implementation reduces the number of simultaneous changes.

\paragraph{Phase A: geometry and timestamp front end.}
Define CoG, IMU, and GNSS antenna points; add a common GNSS sample type with measurement
time, local-NED position, velocity, covariance, validity, and age.  Verify antenna
lever-arm correction independently using truth attitude/gyro data.

\paragraph{Phase B: OU--III joint navigation block.}
OU--III is the best first target because its $S_w=0$ regularizer already leaves
$\vct p_w$ and $\vct v_w$ available for direct GNSS sum measurements.  Add
$(\vct r_n,\vct u_n,\vct a_n)$, cross-covariance, and event-time GNSS updates while
reusing the existing IMU CoG lever-arm model.

\paragraph{Phase C: OU--II parity.}
Add the same navigation block but retain zero-$p_w$/zero-$v_w$ virtual measurements.
This creates a useful experiment on whether GNSS and explicit navigation reduce the
need for those stronger wave anchors.

\paragraph{Phase D: TFG geometry.}
Implement the GNSS residual in right-invariant coordinates, first with an additive
navigation block and then, if justified, with navigation as additional two-frame world
columns.

\paragraph{Phase E: NLO parity.}
Add timestamped 3-D position/velocity input and antenna-to-CoG correction.  Then
re-derive the translational observer if a simultaneous absolute-trajectory and
bounded-wave output is required.

\section{Validation Methodology}
\label{sec:validation}

No result should be reported from a single deterministic realization.  A minimum
validation bundle should use paired wave/noise seeds and truth at the CoG, IMU, and
antenna points.  Every scenario should log raw GNSS observations, corrected CoG
observations, filter innovations, aiding mode, lever-arm contributions, navigation
states, wave states, attitude/bias states, and covariance or observer confidence.

Primary navigation metrics should include horizontal and vertical position RMS,
velocity RMS, outage drift rate, reacquisition overshoot, and time to recover nominal
error.  Primary wave metrics should include per-axis displacement RMS, vertical error
as percent of $H_s$, spectral gain/phase in the wave band, and leakage between
trajectory and wave components.  The split should be scored directly:
\begin{align}
E_{n\rightarrow w}
&=\frac{\|\mathcal P_w(\hat{\vct p}_w-\vct p_w^{\rm ref})\|}
        {\|\mathcal P_n\vct r_n^{\rm ref}\|+\epsilon},\\
E_{w\rightarrow n}
&=\frac{\|\mathcal P_n(\hat{\vct r}_n-\vct r_n^{\rm ref})\|}
        {\|\mathcal P_w\vct p_w^{\rm ref}\|+\epsilon},
\label{eq:leakage-metrics}
\end{align}
where $\mathcal P_n$ and $\mathcal P_w$ are declared slow and wave-band projections
used only for scoring, not for constructing the estimate.

Consistency metrics for the Kalman families should include GNSS position/velocity NIS
and covariance coverage.  NLO should report normalized innovation energy and observer
transient envelopes instead of pretending to have a Kalman covariance.

\section{Required Study Matrix}
\label{sec:studies}

The following studies are required before turning this design paper into a performance
paper.

\subsection{Reference-point and lever-arm studies}
Use zero wave motion with commanded roll, pitch, yaw, and angular acceleration so the
truth consists almost entirely of lever-arm motion.  Sweep IMU and antenna offsets,
including zero, realistic deck offsets, and deliberately exaggerated baselines.  Test
survey error in magnitude and direction.  For the IMU correction, sweep angular-
acceleration smoothing and compare uncompensated, truth-compensated, and estimated
compensation.

\subsection{Absolute-navigation studies}
Test stationary, constant-velocity, constant-turn-rate, acceleration/deceleration,
and curved trajectories with no waves.  These runs establish whether the slow
navigation model tracks real vessel motion without leaking it into $\vct p_w$.
Include current-driven translation that changes velocity without a corresponding wind
or heading change.

\subsection{Wave-plus-trajectory studies}
Superimpose JONSWAP and PM--Stokes wave realizations on straight sailing, turns, and
slow acceleration.  Include head, following, beam, and quartering seas plus crossing
seas.  The key metric is not only total CoG error but whether
$\hat{\vct r}_n+\hat{\vct p}_w$ is correct \emph{and} each component matches its truth.

\subsection{GNSS-rate, latency, and asynchrony studies}
Sweep 1, 5, 10, and 20~Hz aiding; deterministic and jittered latency; and clock offsets
of 0, 50, 100, and 250~ms.  Compare exact rewind/replay with forward-propagated delayed
measurements.  Include position-only, velocity-only, and combined aiding.

\subsection{Outage and reacquisition studies}
Inject outages of 1, 5, 30, 120, and 600~s at different wave phases and maneuver
states.  Reacquire with clean fixes, biased fixes, and isolated outliers.  Compare hard
acceptance, NIS-gated soft start, and velocity-first reacquisition.  Verify that wave
phase remains continuous when absolute navigation is corrected.

\subsection{GNSS quality and vertical-aiding studies}
Sweep realistic horizontal/vertical covariance, multipath bursts, degraded satellite
geometry, and vertical-bias offsets.  Compare 2-D position plus 3-D velocity against
full 3-D position/velocity.  Determine when low-quality GNSS height improves or harms
heave.

\subsection{Architecture ablations}
For OU--III, compare joint navigation, cascaded navigation, and no GNSS.  For OU--II,
compare fixed zero anchors with GNSS-aware anchor weakening.  For TFG, compare additive
navigation augmentation with group-column augmentation.  For NLO, compare current
horizontal GNSS semantics, lever-arm-corrected GNSS, outage virtual aiding, and the
re-derived trajectory/wave split.  All comparisons should use the same sensor data and
reference trajectories.

\section{Results Reporting Contract}
\label{sec:results}

This manuscript presently contains no GNSS-aided performance table because the joint
trajectory/wave estimator and the corresponding cross-family study have not yet been
implemented.  A future Results section should report, at minimum:
\begin{itemize}
  \item absolute CoG position and velocity errors in aided, outage, and reacquisition
  intervals;
  \item wave surge/sway/heave error and wave-band gain/phase;
  \item navigation-to-wave and wave-to-navigation leakage;
  \item attitude, gyro-bias, and accelerometer-bias changes attributable to GNSS;
  \item innovation consistency and rejection statistics;
  \item lever-arm ablations and survey-error sensitivity;
  \item latency and GNSS-rate sensitivity;
  \item filter-family comparison using identical measurement streams.
\end{itemize}
Any result that reports only total antenna or CoG position error without the separated
wave and trajectory components is insufficient to validate the central claim of this
paper.

\section{Discussion}
\label{sec:discussion}

The most important design decision is semantic rather than numerical.  A marine wave
filter cannot become an absolute navigator merely by sending latitude/longitude to an
existing position pseudo-measurement.  If that position state is explicitly
zero-mean, the measurement models are contradictory.  The solution is to introduce a
slow navigation trajectory and let physical sensors observe the sum of trajectory and
wave motion.

The same reasoning clarifies lever arms.  There are two different baselines and two
different kinematic effects.  The IMU-to-CoG offset modifies measured specific force
through angular acceleration and centripetal acceleration.  The GNSS-antenna-to-CoG
offset modifies measured position through attitude and measured velocity through
angular rate.  Using only an IMU-to-GNSS baseline can hide which reference point the
state represents; expressing both offsets from CoG makes the model auditable.

GNSS aiding should substantially improve horizontal low-frequency observability, which
is a known weakness of unaided marine inertial motion estimation.  But the study must
verify that this benefit does not come from attenuating the desired wave signal.  The
joint model provides a mechanism for avoiding that trade: GNSS constrains the slow
trajectory and the physical sum, while the wave prior and virtual measurements shape
only the bounded component.

\section{Conclusion}
\label{sec:conclusion}

A coherent GNSS-aided marine inertial estimator needs three things at once: a declared
reference point, correct rigid-body sensor geometry, and separate semantics for
absolute trajectory and wave motion.  The proposed model writes
\[
\vct r_C=\vct r_n+\vct p_w,\qquad
\vct v_C=\vct u_n+\vct v_w,\qquad
\vct a_C=\vct a_n+\vct a_w,
\]
and lets the GNSS antenna observe those sums through its attitude-dependent lever arm.
The IMU observes the same CoG acceleration through its own rotational lever-arm term.
OU--II and OU--III already contain the latter correction and the covariance-update
machinery required for position/velocity aiding, but absolute GNSS should be added by
augmenting the state rather than reusing a zero-mean wave state as global position.
TFG requires a group-consistent measurement extension, while NLO requires timestamped
lever-arm-corrected satellite aiding and a re-derived state split if it is to output
absolute navigation and bounded wave motion simultaneously.  The resulting framework
also gives clean outage behavior: absolute navigation uncertainty may grow while the
wave estimator remains drift bounded, and reacquisition can correct the trajectory
without resetting wave phase.  The next step is implementation followed by the paired
lever-arm, latency, outage, maneuver, and directional-sea studies specified here.

\bibliographystyle{IEEEtran}
\bibliography{w3d,w3d-baselines}
\end{document}
